\chapter{Conclusion}\label{ch:conclusion}

This thesis set out to study Predictive Coding from scratch and to answer three questions about
it honestly. The answers, taken together, form a sober but useful picture.

To \textbf{RQ1}---whether PC outperforms backpropagation under a fair comparison---the answer is
\emph{no, not at this scale}. Across bulk accuracy, noise robustness, sample efficiency,
domain-incremental and class-incremental learning, a loss-matched and per-method-tuned PC network
is statistically indistinguishable from a matched backpropagation baseline
(\Cref{sec:exp-pcbp}). The most faithful test we could construct---an architecture-exact
replication of \citet{song2024prospective}'s alternating Split-FashionMNIST protocol, with the
training budget treated as a controlled axis and ten seeds---agrees with this parity verdict
rather than overturning it: no PC-versus-BP difference exceeds one standard deviation at any
budget (\Cref{sec:exp-song}). The contribution here is not the null itself but its explanation.
We isolate three confounds---the loss function, the stability--plasticity trade-off in forgetting
metrics, and an untuned baseline learning rate---each of which manufactures a spurious
``PC advantage'' in a naive comparison, and we offer the resulting checklist as a transferable
guard for the field. That the correct interpretation of the replication itself only emerged after
an adversarial three-lens review, which caught us misattributing a learning-rate-robustness claim
to an interference figure, is reported as both a methodological tool and a cautionary tale
(\Cref{ch:discussion}).

To \textbf{RQ2}---where PC's inference cost goes and whether it can be reduced---the answer is
concrete. PC inference at MNIST scale is launch-overhead-bound: the framework issues tens of tiny
kernel launches per settling step and leaves the GPU half-idle (\Cref{sec:exp-kernel}). A fused
CUDA kernel that runs the entire settling loop in a single launch, with the layer states resident
on chip, removes that overhead and is---to our knowledge---the first hand-written settling kernel
for PC. It delivers a $3.2\times$ speedup at small batch, generalises to arbitrary depth, and was
fast enough to serve as the actual backend of our exact-replication study, where it ran
$1.45\times$ faster end-to-end with per-seed-identical results. The same residency that makes it
fast at small batch is architecturally opposed to the register-blocked tiling a cuBLAS-class
matrix multiply needs, so beyond batch $\sim$$2048$ the standard backend wins---a genuine systems
finding, not a missing optimisation.

To \textbf{RQ3}---whether one PC mechanism serves many tasks---the answer is a qualified yes, with
an important precondition. A \emph{discriminatively} trained PC cannot generate; only when the
network is oriented generatively (label$\rightarrow$image) does the single settling rule yield
classification, class-conditional generation, inpainting, and anomaly detection at once
(\Cref{sec:exp-gen}). The ``one mechanism, many tasks'' property is real, but it is a property of
the generative orientation, not of PC in general.

The deliverable of this work is therefore a reproducible, honestly-scoped baseline plus three
transferable insights: the launch-overhead characterisation of PC inference, the PC-versus-BP
confound checklist, and an adversarial-verification protocol that caught a framing error before it
reached the page. None of these is a paradigm shift, and none is presented as one. In a subfield
whose central empirical question---\emph{when}, if ever, PC does something backpropagation does
not---remains open, a clean negative result, a precise systems measurement, and a reusable
methodological guard are, we argue, exactly the kind of contribution that moves the discussion
forward. \Cref{ch:limitations} lays out what would be required to extend these claims: greater
scale and depth, an Equilibrium-Propagation comparison arm, harder out-of-distribution tests, and
the algorithmic and kernel-level work needed to make deep PC both stable and fast. The code,
tests, figures, and artifacts are released so that the next study can start where this one ends.
